\subsection{复合积与扩张模的连接同态}

命题 5. --- 设

\[
0 \to {\bf M} ^ {\prime} \stackrel {f} {\to} {\bf M} \stackrel {g} {\to} {\bf M} ^ {\prime \prime} \to 0\tag{8}
\]

是左A-模的一个正合序列， \(\theta\in\mathrm{Ext}_{\mathbf{A}}^{1}\left(\mathbf{M}^{\prime\prime},\mathbf{M}^{\prime}\right)\) 为相应的类，N为左A-模，n为整数。

\begin{enumerate}
\def\labelenumi{\alph{enumi})}
\item
  连接同态 \(\delta^{n}(\mathbf{N},\mathcal{E}):\operatorname{Ext}_{\mathbf{A}}^{n}(\mathbf{N},\mathbf{M}^{\prime \prime})\to \operatorname{Ext}_{\mathbf{A}}^{n + 1}(\mathbf{N},\mathbf{M}^{\prime})\) 是复合积 \(\alpha \mapsto \theta \circ \alpha\) ，即用 \(\theta\) 作复合而得.
\item
  连接同态 \(\delta^{n}(\mathcal{E},\mathbf{N}):\operatorname{Ext}_{\mathbf{A}}^{n}(\mathbf{M}^{\prime},\mathbf{N})\to\operatorname{Ext}_{\mathbf{A}}^{n+1}(\mathbf{M}^{\prime\prime},\mathbf{N})\) 是复合积 \(\alpha\mapsto(-1)^{n+1}\alpha\circ\theta\) ，即用 \((-1)^{n+1}\theta\) 作复合而得.
\item
  考虑一个交换图
\end{enumerate}

\[
\begin{array}{l} 0 \xrightarrow {} \mathrm{M} ^ {\prime} \xrightarrow {f} \mathrm{M} \xrightarrow {g} \mathrm{M} ^ {\prime \prime} \xrightarrow {} 0 \\ 0 \xrightarrow {} \mathrm{M} ^ {\prime} \xrightarrow {e _ {\mathrm{M} ^ {\prime}}} \mathrm{I} ^ {0} (\mathrm{M} ^ {\prime}) \xrightarrow {\delta^ {0}} \mathrm{I} ^ {1} (\mathrm{M} ^ {\prime}). \end{array}
\]

根据定义， \(\theta\) 是 \(-v^{1}\in\mathrm{Homgr}_{\mathbf{A}}^{1}(\mathbf{M}^{\prime\prime},\mathbf{I}(\mathbf{M}^{\prime}))\) 的类。再设 \(\alpha\in\mathrm{Ext}_{\mathbf{A}}^{n}(\mathbf{N},\mathbf{M}^{\prime\prime})\) ，它由元素 \(a\) （属于 \(\mathrm{Homgr}_{\mathbf{A}}^{n}(\mathbf{L}(\mathbf{N}),\mathbf{M}^{\prime\prime})\) ）来表示。根据构造， \(\delta^{n}(\alpha)\) 如下获得：将 \(a^{n}\in\mathrm{Hom}_{\mathbf{A}}(\mathbf{L}_{n}(\mathbf{N}),\mathbf{M}^{\prime\prime})\) 提升为

\[
b \in \operatorname{Hom} _ {\mathbf {A}} \left(\mathrm{L} _ {n} (\mathrm{N}), \mathrm{M}\right),
\]

而 \(\delta^n (\alpha)\) 是 \(e_{\mathbf{M}'}\circ c\) 的类，其中 \(c\in \mathrm{Hom}_{\mathbf{A}}\left(\mathbf{L}_{n + 1}(\mathbf{N}),\mathbf{M}'\right)\) 满足

\[
f \circ c = \mathrm{D} b = (- 1) ^ {n + 1} b \circ d _ {n + 1}.
\]

因此我们有一个交换图

\[
\begin{array}{c} L _ {n + 1} (N) \xrightarrow {d _ {n + 1}} L _ {n} (N) \\ (- 1) ^ {n + 1} c \Big \downarrow \\ 0 \longrightarrow M ^ {\prime} \xrightarrow {f} M \xrightarrow {g} M ^ {\prime \prime} \longrightarrow 0 \\ 1 _ {M ^ {\prime}} \Big \downarrow \\ 0 \longrightarrow M ^ {\prime} \xrightarrow {e _ {M ^ {\prime}}} I ^ {0} (M ^ {\prime}) \xrightarrow {\delta^ {0}} I ^ {1} (M ^ {\prime}). \end{array}
\]

但在 \(\operatorname{Homgr}_{\mathbf{A}}(\mathbf{L}(\mathbf{N}), \mathbf{l}(\mathbf{M}'))\) 中，我们有

\[
\mathrm{D} \left(v ^ {0} \circ b\right) = \delta^ {0} \circ v ^ {0} \circ b - (- 1) ^ {n} v ^ {0} \circ b \circ d _ {n + 1} = v ^ {1} \circ a ^ {n} + e _ {\mathrm{M} ^ {\prime}} \circ c.
\]

\(e_{\mathbf{M}'} \circ c\) 与 \((-v^1) \circ a\) 的类在 \(\operatorname{Ext}_{\mathbf{A}}^{n+1}(\mathbf{N}, \mathbf{M}')\) 中相等，由此得 \(a\) ).

\begin{enumerate}
\def\labelenumi{\alph{enumi})}
\setcounter{enumi}{1}
\tightlist
\item
  考虑一个交换图
\end{enumerate}

\[
\begin{array}{c} \mathrm{L} _ {1} (\mathrm{M} ^ {\prime \prime}) \xrightarrow {d _ {1}} \mathrm{L} _ {0} (\mathrm{M} ^ {\prime \prime}) \xrightarrow {p _ {\mathrm{M} ^ {\prime \prime}}} \mathrm{M} ^ {\prime \prime} \longrightarrow 0 \\ \Biggl \downarrow u _ {1} \\ 0 \longrightarrow \mathrm{M} ^ {\prime} \xrightarrow {f} \mathrm{M} \xrightarrow {g} \mathrm{M} ^ {\prime \prime} \longrightarrow 0. \end{array}
\]

根据定义， \(\theta\) 是 \(-u_{1}\in\mathrm{Homgr}_{\mathbf{A}}^{1}(\mathrm{L}(\mathbf{M}^{\prime\prime}),\mathbf{M}^{\prime})\) 的类。再设 \(\alpha\in\mathrm{Ext}^{n}(\mathbf{M}^{\prime},\mathbf{N})\) ，它由元素 \(a\) （属于 \(\mathrm{Homgr}_{\mathbf{A}}^{n}(\mathbf{M}^{\prime},\mathrm{I}(\mathbf{N}))\) ）来表示。根据构造， \(\delta^{n}(\alpha)\) 如下获得：将 \(a^{n}\in\mathrm{Hom}_{\mathbf{A}}(\mathbf{M}^{\prime},\mathrm{I}^{n}(\mathbf{N}))\) 扩充为

\[
b \in \operatorname{Hom} _ {\mathbf {A}} (\mathbf {M}, \mathrm{I} ^ {n} (\mathbf {N}))
\]

而 \(\delta^n (\alpha)\) 是 \(c\circ p_{\mathbf{M}''}\) 的类，其中 \(c\in \mathrm{Hom}_{\mathbf{A}}(\mathbf{M}^{\prime \prime},\mathrm{I}^{n + 1}(\mathbf{N}))\) 满足

\[
g \circ c = D b = \delta^ {n + 1} \circ b.
\]

因此我们有一个交换图

\[
\begin{array}{c} L _ {1} (M ^ {\prime \prime}) \xrightarrow {d _ {1}} L _ {0} (M ^ {\prime \prime}) \xrightarrow {p _ {M ^ {\prime \prime}}} M ^ {\prime \prime}. \longrightarrow 0 \\ u _ {1} \Biggl \downarrow \quad \quad \quad \quad \quad \quad \quad \quad \quad \quad \quad \quad \quad \quad \quad \quad \quad \quad \quad \quad \quad \quad \quad \quad \quad \quad \quad \quad \quad \quad \quad \quad \quad \quad \quad \quad \quad \quad \quad \quad \quad \quad \quad \quad \quad \quad \quad 1 _ {M ^ {\prime}} \\ 0 \longrightarrow M ^ {\prime} \xrightarrow {f} M \xrightarrow {g} M ^ {\prime \prime} \longrightarrow 0 \\ a ^ {n} \Biggl \downarrow \quad \quad \quad \quad \quad \quad \quad \quad \quad \quad \quad \quad c \Biggl \downarrow \\ I ^ {n} (N) \xrightarrow {\delta^ {n}} I ^ {n + 1} (N). \end{array}
\]

但在 \(\operatorname{Homgr}_{\mathbf{A}}(\mathbf{L}(\mathbf{M}^{\prime \prime}),\mathbf{l}(\mathbf{N}))\) 中，我们有

\[
\mathrm{D} (b \circ u _ {0}) = \delta^ {n} \circ b \circ u _ {0} - (- 1) ^ {n} b \circ u _ {0} \circ c = c \circ p - (- 1) ^ {n} a ^ {n} \circ u _ {1}.
\]

\(c \circ p\) 与 \((-1)^{n+1} a^n \circ (-u_1)\) 的类在 \(\operatorname{Ext}_{\mathbf{A}}^{n+1} (\mathbf{M}''\) , \(\mathbf{N}\) ) 中因此相等，由此得 \(b\) ).

推论 1. --- a) 连接同态 Hom \(_{A}\) (M'', M'') → Ext \(_{A}^{1}\) (M'', M') 将 1 \(_{M''}\) 映至 θ。

\begin{enumerate}
\def\labelenumi{\alph{enumi})}
\setcounter{enumi}{1}
\tightlist
\item
  连接同态 \(\operatorname{Hom}_{\mathbf{A}}(\mathbf{M}^{\prime},\mathbf{M}^{\prime})\to\operatorname{Ext}_{\mathbf{A}}^{1}(\mathbf{M}^{\prime\prime},\mathbf{M}^{\prime})\) 将 \(1_{M'}\) 映至 \(-\theta\).
\end{enumerate}

推论 2. --- 考虑左 A-模的两个正合序列

\[
0 \rightarrow \mathbf {M} ^ {\prime} \rightarrow \mathbf {M} \rightarrow \mathbf {M} ^ {\prime \prime} \rightarrow 0
\]

\[
0 \rightarrow \mathrm{N} ^ {\prime} \rightarrow \mathrm{N} \rightarrow \mathrm{N} ^ {\prime \prime} \rightarrow 0.
\]

则连接同态的复合同态

\[
\operatorname{Ext} _ {\mathbf {A}} ^ {n} \left(\mathrm{M} ^ {\prime}, \mathrm{N} ^ {\prime \prime}\right)\rightarrow \operatorname{Ext} _ {\mathbf {A}} ^ {n + 1} \left(\mathrm{M} ^ {\prime \prime}, \mathrm{N} ^ {\prime \prime}\right)\rightarrow \operatorname{Ext} _ {\mathbf {A}} ^ {n + 2} \left(\mathrm{M} ^ {\prime \prime}, \mathrm{N} ^ {\prime}\right)
\]

与

\[
\operatorname{Ext} _ {\mathbf {A}} ^ {n} \left(\mathrm{M} ^ {\prime}, \mathrm{N} ^ {\prime \prime}\right)\rightarrow \operatorname{Ext} _ {\mathbf {A}} ^ {n + 1} \left(\mathrm{M} ^ {\prime}, \mathrm{N} ^ {\prime}\right)\rightarrow \operatorname{Ext} _ {\mathbf {A}} ^ {n + 2} \left(\mathrm{M} ^ {\prime \prime}, \mathrm{N} ^ {\prime}\right)
\]

是相反的。

事实上，若 \(\theta_{1}\) , \(\theta_{2}\) 是与给定正合序列相关联的类，且若 \(\alpha \in \operatorname{Ext}_{\mathbf{A}}^{n}(\mathbf{M}', \mathbf{M}'')\) ，则 \(\alpha\) 的像分别为

\[
\theta_ {2} \circ ((- 1) ^ {n + 1} \alpha \circ \theta_ {1}) \quad \text { and } \quad (\theta_ {2} \circ \alpha) \circ ((- 1) ^ {n + 2} \theta_ {1}).
\]

考虑左 A-模的一个正合序列

\[
0 \rightarrow \mathrm{N} \rightarrow \mathrm{R} _ {n} \xrightarrow {f _ {n}} \mathrm{R} _ {n - 1} \xrightarrow {f _ {n - 1}} \dots \rightarrow \mathrm{R} _ {1} \xrightarrow {f _ {1}} \mathrm{M} \rightarrow 0\tag{S}
\]

并设 \(K_{0} = M\) , \(K_{i} = Ker f_{i}, i = 1, \ldots, n - 1\) , \(K_{n} = N\) 。因此我们有正合序列

\[
0 \to \mathbf {K} _ {i} \to \mathbf {R} _ {i} \to \mathbf {K} _ {i - 1} \to 0, \quad 1 \leqslant i \leqslant n,\tag{9}
\]

与之相关联的，对每个左 A-模 P，有连接同态

\[
\begin{array}{r l}&{\mathrm{Ext} _ {\mathbf {A}} ^ {m} \left(\mathrm{P}, \mathbf {K} _ {i - 1}\right)\rightarrow \mathrm{Ext} _ {\mathbf {A}} ^ {m + 1} \left(\mathrm{P}, \mathbf {K} _ {i}\right),}\\&{\qquad \mathrm{Ext} _ {\mathbf {A}} ^ {m} \left(\mathbf {K} _ {i}, \mathrm{P}\right)\rightarrow \mathrm{Ext} _ {\mathbf {A}} ^ {m + 1} \left(\mathbf {K} _ {i - 1}, \mathrm{P}\right),}\end{array}
\]

从而通过迭代连接同态的复合，与 (S) 相关联

\[
\begin{array}{l} \delta^ {m} (\mathrm{P}, \mathcal {S}): \operatorname{Ext} _ {\mathrm{A}} ^ {m} (\mathrm{P}, \mathrm{M}) \to \operatorname{Ext} _ {\mathrm{A}} ^ {m + n} (\mathrm{P}, \mathrm{N}) \\ \delta^ {m} (\mathcal {S}, \mathrm{P}): \operatorname{Ext} _ {\mathrm{A}} ^ {m} (\mathrm{N}, \mathrm{P}) \to \operatorname{Ext} _ {\mathrm{A}} ^ {m + n} (\mathrm{M}, \mathrm{P}). \end{array}
\]

推论 3. --- 如果 \(\theta\in\mathrm{Ext}_{\mathbf{A}}^{n}(\mathbf{M},\mathbf{N})\) 是与正合序列 \((\mathcal{S})\) 相关联的类，则我们有

\[
\delta^ {m} (\mathrm{P}, \mathscr {S}) (\alpha) = \theta \circ \alpha , \quad \delta^ {m} (\mathscr {S}, \mathrm{P}) (\beta) = (- 1) ^ {m n + n (n + 1) / 2} \beta \circ \theta .
\]

若 \(\theta_{i}\in\mathrm{Ext}_{\mathbf{A}}^{1}(\mathbf{K}_{i-1},\mathbf{K}_{i})\) 是与正合序列 (9) 相关联的类，根据命题 5，我们有

\[
\begin{array}{l} \delta^ {m} (\mathrm{P}, \mathcal {S}) (\alpha) = \theta_ {n} \circ \dots \circ \theta_ {2} \circ \theta_ {1} \circ \alpha \\ \delta^ {m} (\mathcal {S}, \mathrm{P}) (\beta) = (- 1) ^ {(m + 1) + \dots + (m + n)} \beta \circ \theta_ {n} \circ \dots \circ \theta_ {1}. \end{array}
\]

此外，根据命题 3（X，第 118 页），我们有 \(\theta = \theta_{n} \circ \cdots \circ \theta_{1}\) 。推论立即由此及关系式（E，III，第44页）得出

\[
(m + 1) + \dots + (m + n) = m n + n (n + 1) / 2.
\]

推论4. --- 若每个模 \(\mathbf{R}_i\) , \(i = 1, \ldots, n\) 都是内射的（相应地，投射的），则映射 \(\alpha \mapsto \theta \circ \alpha\) （相应地， \(\alpha \mapsto \alpha \circ \theta\) ）从 \(\operatorname{Ext}_{\mathbf{A}}^{m}(\mathbf{P}, \mathbf{M})\) 到 \(\operatorname{Ext}_{\mathbf{A}}^{m+n}(\mathbf{P}, \mathbf{N})\) （相应地，从 \(\operatorname{Ext}_{\mathbf{A}}^{m}(\mathbf{N}, \mathbf{P})\) 到 \(\operatorname{Ext}_{\mathbf{A}}^{m+n}(\mathbf{M}, \mathbf{P})\) ）对每个 A-模 P 和每个整数 \(m > 0\) 都是双射的。这确实由推论 3 及正合序列

\[
\operatorname{Ext} _ {\mathbf {A}} ^ {m + i - 1} \left(\mathrm{P}, \mathrm{R} _ {i}\right)\rightarrow \operatorname{Ext} _ {\mathbf {A}} ^ {m + i - 1} \left(\mathrm{P}, \mathrm{K} _ {i - 1}\right)\rightarrow \operatorname{Ext} _ {\mathbf {A}} ^ {m + i} \left(\mathrm{P}, \mathrm{K} _ {i}\right)\rightarrow \operatorname{Ext} _ {\mathbf {A}} ^ {m + i} \left(\mathrm{P}, \mathrm{R} _ {i}\right)
\]

\[
\left(\text { resp. } \operatorname{Ext} _ {\mathbf {A}} ^ {m + i - 1} \left(\mathrm{R} _ {i}, \mathrm{P}\right)\rightarrow \operatorname{Ext} _ {\mathbf {A}} ^ {m + i - 1} \left(\mathrm{K} _ {i}, \mathrm{P}\right)\rightarrow \operatorname{Ext} _ {\mathbf {A}} ^ {m + i} \left(\mathrm{K} _ {i - 1}, \mathrm{P}\right)\rightarrow \operatorname{Ext} _ {\mathbf {A}} ^ {m + i} \left(\mathrm{R} _ {i}, \mathrm{P}\right)\right),
\]

得出，其两端项按假设为零。

注记. --- 第3至6节的定义和命题适用于右A-模，它们被视为环A的反环A°上的左模。

